% =====================================================
% 03_theory.tex — ~1.5 pg
% =====================================================
\todo{draft — full content in theory/*.md; this is the LaTeX-ified short version}

\subsection{Setup}
Let $\W \cong \R^{LD}$ be the StyleGAN $W^+$ latent space and
$\Img = [-1,1]^{3 \times H \times W}$ the image domain. The generator
$G : \W \to \Img$ is $C^2$ at almost every $wp$ (true for the
ReLU/AdaIN architectures we consider, off a measure-zero set).
The \emph{image manifold} is $\Mfd = G(\W) \subset \Img$.

\subsection{Pushforward and first-order saliency}
The differential $\dG{wp} : T_{wp} \W \to T_{G(wp)} \Mfd$ sends a
latent tangent vector to its first-order image response. For a unit
direction $b \in T \W$,
\begin{equation}
  \dG{wp}(b) \;=\; \frac{\partial G(wp + \alpha b)}{\partial \alpha}\bigg|_{\alpha = 0}
\end{equation}
is exactly what \texttt{torch.func.jvp(G, (wp,), (b,))} returns. We
define the \emph{first-order saliency map}
\begin{equation}
  S^{(1)}_{h,w}(wp, b) \;=\; \tfrac{1}{3}\sum_c \big| [\dG{wp}(b)]_{c,h,w} \big|.
\end{equation}
Its prior-averaged version $\bar{S}^{(1)}(b)$ is our pixel-attribution
counterpart to Grad-CAM~\cite{Selvaraju2017GradCAM}.

Forward mode is the right algorithmic choice: the perturbation lives
in $\R^1$ (one scalar magnitude) while the output lives in
$\R^{3HW} \approx 2\!\times\!10^5$. Reverse mode would require $3HW$
backward passes for the same per-pixel information.

\subsection{Second-order pushforward and extrinsic curvature}
The Hessian-vector product
\begin{equation}
  H(wp, b) \;=\; \frac{\partial^2 G(wp + \alpha b)}{\partial \alpha^2}\bigg|_{\alpha=0}
  \;\in\; \R^{3HW}
\end{equation}
is computable in one \emph{composed} \jvp: $\jvp \circ \jvp$. Its
projection onto the normal bundle of $\Mfd$ is the second fundamental
form of the embedding $\Mfd \hookrightarrow \R^{3HW}$. The magnitude
$|H|$ upper-bounds the magnitude of the intrinsic extrinsic curvature
and is in any case a directly meaningful "non-linearity per pixel"
quantity. The \emph{non-linearity ratio}
\begin{equation}
  \rho_{h,w}(wp, b) \;=\; \frac{S^{(2)}_{h,w}(wp, b)}{S^{(1)}_{h,w}(wp, b) + \epsilon}
\end{equation}
distinguishes directions whose edits remain in a linear regime
($\rho \ll 1$) from those that introduce topological transitions
($\rho \gg 1$).

\subsection{The six claims}
\begin{claim}[C1: measurable curvature]
For every smooth direction $b$ and base $wp$, $H(wp, b)$ is
well-defined and computable by composed forward \jvp.
\end{claim}

\begin{claim}[C2: curvature predicts topology change]
Directions with $\bar{\rho} \gg 1$ produce edits that change the set
of image-space level sets (object insertion / removal), whereas
$\bar{\rho} \approx 1$ directions only repaint pixel values.
\end{claim}

\begin{claim}[C3: boundaries are layer-localised sections]
A boundary direction $b \in \R^D$ is a section of the layer fibration
$\pi : \W \to \{1,\dots,L\}$. Applied off its canonical fibre, $b$
produces large but semantically incoherent pushforwards.
\end{claim}

\begin{claim}[C4: compositional failure follows from the mixed Hessian]
Compositional editing failure between directions $b_a, b_b$ is
predicted by $\E_{wp}\|d^2 G_{wp}(b_a, b_b)\|$, normalised by
$\E\|\dG{wp} b_a\| \cdot \E\|\dG{wp} b_b\|$.
\end{claim}

\begin{claim}[C5: encoders are charts]
An encoder $E : \Img \to \W$ approximates a coordinate chart on
$\Mfd$. Its quality is best measured by chart consistency
(agreement of derivatives) rather than reconstruction error.
\end{claim}

\begin{claim}[C6: random directions stratify the manifold]
K-means clustering of $\Phi(b) = \E_{wp} \bar{S}^{(1)}(wp, b)$ for
$N \gtrsim 128$ random unit directions produces a stratification of
$\Mfd$ whose largest clusters recover hand-curated attribute boundaries
without supervision.
\end{claim}

Each claim is independently testable. Several are causally connected
— C4 is a consequence of the Hessian formula and C1, and C6 is the
many-direction extension of C1+C2. We present quantitative validation
for all six in \S\ref{sec:experiments}.
